% Options for packages loaded elsewhere
\PassOptionsToPackage{unicode}{hyperref}
\PassOptionsToPackage{hyphens}{url}
\documentclass[
]{article}
\usepackage{xcolor}
\usepackage[margin=1in]{geometry}
\usepackage{amsmath,amssymb}
\setcounter{secnumdepth}{-\maxdimen} % remove section numbering
\usepackage{iftex}
\ifPDFTeX
  \usepackage[T1]{fontenc}
  \usepackage[utf8]{inputenc}
  \usepackage{textcomp} % provide euro and other symbols
\else % if luatex or xetex
  \usepackage{unicode-math} % this also loads fontspec
  \defaultfontfeatures{Scale=MatchLowercase}
  \defaultfontfeatures[\rmfamily]{Ligatures=TeX,Scale=1}
\fi
\usepackage{lmodern}
\ifPDFTeX\else
  % xetex/luatex font selection
\fi
% Use upquote if available, for straight quotes in verbatim environments
\IfFileExists{upquote.sty}{\usepackage{upquote}}{}
\IfFileExists{microtype.sty}{% use microtype if available
  \usepackage[]{microtype}
  \UseMicrotypeSet[protrusion]{basicmath} % disable protrusion for tt fonts
}{}
\makeatletter
\@ifundefined{KOMAClassName}{% if non-KOMA class
  \IfFileExists{parskip.sty}{%
    \usepackage{parskip}
  }{% else
    \setlength{\parindent}{0pt}
    \setlength{\parskip}{6pt plus 2pt minus 1pt}}
}{% if KOMA class
  \KOMAoptions{parskip=half}}
\makeatother
\usepackage{color}
\usepackage{fancyvrb}
\newcommand{\VerbBar}{|}
\newcommand{\VERB}{\Verb[commandchars=\\\{\}]}
\DefineVerbatimEnvironment{Highlighting}{Verbatim}{commandchars=\\\{\}}
% Add ',fontsize=\small' for more characters per line
\usepackage{framed}
\definecolor{shadecolor}{RGB}{248,248,248}
\newenvironment{Shaded}{\begin{snugshade}}{\end{snugshade}}
\newcommand{\AlertTok}[1]{\textcolor[rgb]{0.94,0.16,0.16}{#1}}
\newcommand{\AnnotationTok}[1]{\textcolor[rgb]{0.56,0.35,0.01}{\textbf{\textit{#1}}}}
\newcommand{\AttributeTok}[1]{\textcolor[rgb]{0.13,0.29,0.53}{#1}}
\newcommand{\BaseNTok}[1]{\textcolor[rgb]{0.00,0.00,0.81}{#1}}
\newcommand{\BuiltInTok}[1]{#1}
\newcommand{\CharTok}[1]{\textcolor[rgb]{0.31,0.60,0.02}{#1}}
\newcommand{\CommentTok}[1]{\textcolor[rgb]{0.56,0.35,0.01}{\textit{#1}}}
\newcommand{\CommentVarTok}[1]{\textcolor[rgb]{0.56,0.35,0.01}{\textbf{\textit{#1}}}}
\newcommand{\ConstantTok}[1]{\textcolor[rgb]{0.56,0.35,0.01}{#1}}
\newcommand{\ControlFlowTok}[1]{\textcolor[rgb]{0.13,0.29,0.53}{\textbf{#1}}}
\newcommand{\DataTypeTok}[1]{\textcolor[rgb]{0.13,0.29,0.53}{#1}}
\newcommand{\DecValTok}[1]{\textcolor[rgb]{0.00,0.00,0.81}{#1}}
\newcommand{\DocumentationTok}[1]{\textcolor[rgb]{0.56,0.35,0.01}{\textbf{\textit{#1}}}}
\newcommand{\ErrorTok}[1]{\textcolor[rgb]{0.64,0.00,0.00}{\textbf{#1}}}
\newcommand{\ExtensionTok}[1]{#1}
\newcommand{\FloatTok}[1]{\textcolor[rgb]{0.00,0.00,0.81}{#1}}
\newcommand{\FunctionTok}[1]{\textcolor[rgb]{0.13,0.29,0.53}{\textbf{#1}}}
\newcommand{\ImportTok}[1]{#1}
\newcommand{\InformationTok}[1]{\textcolor[rgb]{0.56,0.35,0.01}{\textbf{\textit{#1}}}}
\newcommand{\KeywordTok}[1]{\textcolor[rgb]{0.13,0.29,0.53}{\textbf{#1}}}
\newcommand{\NormalTok}[1]{#1}
\newcommand{\OperatorTok}[1]{\textcolor[rgb]{0.81,0.36,0.00}{\textbf{#1}}}
\newcommand{\OtherTok}[1]{\textcolor[rgb]{0.56,0.35,0.01}{#1}}
\newcommand{\PreprocessorTok}[1]{\textcolor[rgb]{0.56,0.35,0.01}{\textit{#1}}}
\newcommand{\RegionMarkerTok}[1]{#1}
\newcommand{\SpecialCharTok}[1]{\textcolor[rgb]{0.81,0.36,0.00}{\textbf{#1}}}
\newcommand{\SpecialStringTok}[1]{\textcolor[rgb]{0.31,0.60,0.02}{#1}}
\newcommand{\StringTok}[1]{\textcolor[rgb]{0.31,0.60,0.02}{#1}}
\newcommand{\VariableTok}[1]{\textcolor[rgb]{0.00,0.00,0.00}{#1}}
\newcommand{\VerbatimStringTok}[1]{\textcolor[rgb]{0.31,0.60,0.02}{#1}}
\newcommand{\WarningTok}[1]{\textcolor[rgb]{0.56,0.35,0.01}{\textbf{\textit{#1}}}}
\usepackage{graphicx}
\makeatletter
\newsavebox\pandoc@box
\newcommand*\pandocbounded[1]{% scales image to fit in text height/width
  \sbox\pandoc@box{#1}%
  \Gscale@div\@tempa{\textheight}{\dimexpr\ht\pandoc@box+\dp\pandoc@box\relax}%
  \Gscale@div\@tempb{\linewidth}{\wd\pandoc@box}%
  \ifdim\@tempb\p@<\@tempa\p@\let\@tempa\@tempb\fi% select the smaller of both
  \ifdim\@tempa\p@<\p@\scalebox{\@tempa}{\usebox\pandoc@box}%
  \else\usebox{\pandoc@box}%
  \fi%
}
% Set default figure placement to htbp
\def\fps@figure{htbp}
\makeatother
% definitions for citeproc citations
\NewDocumentCommand\citeproctext{}{}
\NewDocumentCommand\citeproc{mm}{%
  \begingroup\def\citeproctext{#2}\cite{#1}\endgroup}
\makeatletter
 % allow citations to break across lines
 \let\@cite@ofmt\@firstofone
 % avoid brackets around text for \cite:
 \def\@biblabel#1{}
 \def\@cite#1#2{{#1\if@tempswa , #2\fi}}
\makeatother
\newlength{\cslhangindent}
\setlength{\cslhangindent}{1.5em}
\newlength{\csllabelwidth}
\setlength{\csllabelwidth}{3em}
\newenvironment{CSLReferences}[2] % #1 hanging-indent, #2 entry-spacing
 {\begin{list}{}{%
  \setlength{\itemindent}{0pt}
  \setlength{\leftmargin}{0pt}
  \setlength{\parsep}{0pt}
  % turn on hanging indent if param 1 is 1
  \ifodd #1
   \setlength{\leftmargin}{\cslhangindent}
   \setlength{\itemindent}{-1\cslhangindent}
  \fi
  % set entry spacing
  \setlength{\itemsep}{#2\baselineskip}}}
 {\end{list}}
\usepackage{calc}
\newcommand{\CSLBlock}[1]{\hfill\break\parbox[t]{\linewidth}{\strut\ignorespaces#1\strut}}
\newcommand{\CSLLeftMargin}[1]{\parbox[t]{\csllabelwidth}{\strut#1\strut}}
\newcommand{\CSLRightInline}[1]{\parbox[t]{\linewidth - \csllabelwidth}{\strut#1\strut}}
\newcommand{\CSLIndent}[1]{\hspace{\cslhangindent}#1}
\setlength{\emergencystretch}{3em} % prevent overfull lines
\providecommand{\tightlist}{%
  \setlength{\itemsep}{0pt}\setlength{\parskip}{0pt}}
\usepackage{bookmark}
\IfFileExists{xurl.sty}{\usepackage{xurl}}{} % add URL line breaks if available
\urlstyle{same}
\hypersetup{
  pdftitle={Template},
  pdfauthor={Jakob, Ben, and David},
  hidelinks,
  pdfcreator={LaTeX via pandoc}}

\title{Template}
\author{Jakob, Ben, and David}
\date{2026-06-09}

\begin{document}
\maketitle

\begin{Shaded}
\begin{Highlighting}[]
\FunctionTok{library}\NormalTok{(tidyverse)}
\FunctionTok{library}\NormalTok{(maps)}

\CommentTok{\# Load data}
\NormalTok{gdp\_data           }\OtherTok{\textless{}{-}} \FunctionTok{read\_csv}\NormalTok{(}\StringTok{"data/raw\_gdp.csv"}\NormalTok{)}
\NormalTok{unemployment\_sex   }\OtherTok{\textless{}{-}} \FunctionTok{read\_csv}\NormalTok{(}\StringTok{"data/raw\_unemployment\_sex.csv"}\NormalTok{)}
\NormalTok{unemployment\_youth }\OtherTok{\textless{}{-}} \FunctionTok{read\_csv}\NormalTok{(}\StringTok{"data/raw\_unemployment\_youth.csv"}\NormalTok{)}

\CommentTok{\# Clean data}
\NormalTok{gdp\_clean }\OtherTok{\textless{}{-}}\NormalTok{ gdp\_data }\SpecialCharTok{\%\textgreater{}\%} 
  \FunctionTok{filter}\NormalTok{(geo }\SpecialCharTok{==} \StringTok{"Netherlands"}\NormalTok{) }\SpecialCharTok{\%\textgreater{}\%}
  \FunctionTok{select}\NormalTok{(TIME\_PERIOD, OBS\_VALUE) }\SpecialCharTok{\%\textgreater{}\%}
  \FunctionTok{rename}\NormalTok{(}\AttributeTok{Year =}\NormalTok{ TIME\_PERIOD, }\AttributeTok{GDP\_Growth =}\NormalTok{ OBS\_VALUE) }\SpecialCharTok{\%\textgreater{}\%}
  \FunctionTok{filter}\NormalTok{(}\SpecialCharTok{!}\FunctionTok{is.na}\NormalTok{(GDP\_Growth))}

\NormalTok{unemp\_clean }\OtherTok{\textless{}{-}}\NormalTok{ unemployment\_sex }\SpecialCharTok{\%\textgreater{}\%} 
  \FunctionTok{filter}\NormalTok{(geo }\SpecialCharTok{==} \StringTok{"Netherlands"}\NormalTok{) }\SpecialCharTok{\%\textgreater{}\%}
  \FunctionTok{select}\NormalTok{(TIME\_PERIOD, OBS\_VALUE) }\SpecialCharTok{\%\textgreater{}\%}
  \FunctionTok{rename}\NormalTok{(}\AttributeTok{Year =}\NormalTok{ TIME\_PERIOD, }\AttributeTok{Unemp\_Rate =}\NormalTok{ OBS\_VALUE) }\SpecialCharTok{\%\textgreater{}\%}
  \FunctionTok{mutate}\NormalTok{(}\AttributeTok{Unemp\_Rate =} \FunctionTok{ifelse}\NormalTok{(Unemp\_Rate }\SpecialCharTok{==} \SpecialCharTok{{-}}\DecValTok{99}\NormalTok{, }\ConstantTok{NA}\NormalTok{, Unemp\_Rate)) }\SpecialCharTok{\%\textgreater{}\%}
  \FunctionTok{filter}\NormalTok{(}\SpecialCharTok{!}\FunctionTok{is.na}\NormalTok{(Unemp\_Rate))}

\NormalTok{youth\_clean }\OtherTok{\textless{}{-}}\NormalTok{ unemployment\_youth }\SpecialCharTok{\%\textgreater{}\%} 
  \FunctionTok{filter}\NormalTok{(geo }\SpecialCharTok{==} \StringTok{"Netherlands"}\NormalTok{) }\SpecialCharTok{\%\textgreater{}\%}
  \FunctionTok{select}\NormalTok{(TIME\_PERIOD, OBS\_VALUE) }\SpecialCharTok{\%\textgreater{}\%}
  \FunctionTok{rename}\NormalTok{(}\AttributeTok{Year =}\NormalTok{ TIME\_PERIOD, }\AttributeTok{Youth\_Unemp\_Rate =}\NormalTok{ OBS\_VALUE) }\SpecialCharTok{\%\textgreater{}\%}
  \FunctionTok{mutate}\NormalTok{(}\AttributeTok{Youth\_Unemp\_Rate =} \FunctionTok{ifelse}\NormalTok{(Youth\_Unemp\_Rate }\SpecialCharTok{==} \SpecialCharTok{{-}}\DecValTok{99}\NormalTok{, }\ConstantTok{NA}\NormalTok{, Youth\_Unemp\_Rate)) }\SpecialCharTok{\%\textgreater{}\%}
  \FunctionTok{filter}\NormalTok{(}\SpecialCharTok{!}\FunctionTok{is.na}\NormalTok{(Youth\_Unemp\_Rate))}

\CommentTok{\# Merge and engineer variables}
\NormalTok{combined\_data }\OtherTok{\textless{}{-}}\NormalTok{ gdp\_clean }\SpecialCharTok{\%\textgreater{}\%}
  \FunctionTok{left\_join}\NormalTok{(unemp\_clean, }\AttributeTok{by =} \StringTok{"Year"}\NormalTok{) }\SpecialCharTok{\%\textgreater{}\%}
  \FunctionTok{left\_join}\NormalTok{(youth\_clean, }\AttributeTok{by =} \StringTok{"Year"}\NormalTok{) }\SpecialCharTok{\%\textgreater{}\%}
  \FunctionTok{mutate}\NormalTok{(}\AttributeTok{Youth\_Vulnerability\_Ratio =}\NormalTok{ Youth\_Unemp\_Rate }\SpecialCharTok{/}\NormalTok{ Unemp\_Rate)}
\end{Highlighting}
\end{Shaded}

\section{Set-up your environment}\label{set-up-your-environment}

\section{Title Page}\label{title-page}

A new sentence Include your names

Include the tutorial group number

Include your tutorial lecturer's name

\section{Part 1 - Identify a Social
Problem}\label{part-1---identify-a-social-problem}

Use APA referencing throughout your document.
\href{https://www.mendeley.com/guides/apa-citation-guide/}{Here's a link
to some explanation.}

\subsection{1.1 Describe the Social
Problem}\label{describe-the-social-problem}

Include the following:

\begin{itemize}
\item
  Why is this relevant? (Smith, 2020)
\item
  \ldots{}
\end{itemize}

\section{Part 2 - Data Sourcing}\label{part-2---data-sourcing}

\subsection{2.1 Load in the data}\label{load-in-the-data}

Preferably from a URL, but if not, make sure to download the data and
store it in a shared location that you can load the data in from. Do not
store the data in a folder you include in the Github repository!

midwest is an example dataset included in the tidyverse package

\subsection{2.2 Provide a short summary of the
dataset(s)}\label{provide-a-short-summary-of-the-datasets}

In this case we see 28 variables, but we miss some information on what
units they are in. We also don't know anything about the year/moment in
which this data has been captured.

\begin{Shaded}
\begin{Highlighting}[]
\NormalTok{inline\_code }\OtherTok{=} \ConstantTok{TRUE}
\end{Highlighting}
\end{Shaded}

These are things that are usually included in the metadata of the
dataset. For your project, you need to provide us with the information
from your metadata that we need to understand your dataset of choice.

\subsection{2.3 Describe the type of variables
included}\label{describe-the-type-of-variables-included}

Think of things like:

\begin{itemize}
\item
  Do the variables contain health information or SES information?
\item
  Have they been measured by interviewing individuals or is the data
  coming from administrative sources?
\end{itemize}

\emph{For the sake of this example, I will continue with the
assignment\ldots{}}

\section{Part 3 - Quantifying}\label{part-3---quantifying}

\subsection{3.1 Data cleaning}\label{data-cleaning}

Our primary empirical data cleaning steps---including filtering
exclusively for the Netherlands territorial boundaries and isolating the
continuous percentage level transformations---were successfully executed
within our master centralized data pipeline at the top of the file to
preserve top-to-bottom reproducibility.

Please use a separate `R block' of code for each type of cleaning. So,
e.g.~one for missing values, a new one for removing unnecessary
variables etc.

\subsection{3.2 Generate necessary
variables}\label{generate-necessary-variables}

Variable 1: Youth Vulnerability Ratio

\begin{Shaded}
\begin{Highlighting}[]
\CommentTok{\# This variable was engineered during our master data pipeline at the top:}
\CommentTok{\# combined\_data \textless{}{-} combined\_data \%\textgreater{}\% mutate(Youth\_Vulnerability\_Ratio = Youth\_Unemp\_Rate / Unemp\_Rate)}
\CommentTok{\# It standardizes the youth unemployment rate as a proportion of aggregate total unemployment.}
\end{Highlighting}
\end{Shaded}

Variable 2

\section{Relational database joins and chronological baseline
alignments}\label{relational-database-joins-and-chronological-baseline-alignments}

\section{were successfully completed using left\_join() functions in the
master
pipeline.}\label{were-successfully-completed-using-left_join-functions-in-the-master-pipeline.}

\subsection{3.3 Visualize temporal
variation}\label{visualize-temporal-variation}

\pandocbounded{\includegraphics[keepaspectratio]{Template_Assignment_files/Template_Assignment_files/figure-latex/temporal_plot-1.pdf}}

\subsection{3.4 Visualize spatial
variation}\label{visualize-spatial-variation}

\pandocbounded{\includegraphics[keepaspectratio]{Template_Assignment_files/Template_Assignment_files/figure-latex/spatial_map-1.pdf}}

Here you provide a description of why the plot above is relevant to your
specific social problem.

\subsection{3.5 Visualize sub-population
variation}\label{visualize-sub-population-variation}

What is the poverty rate by state?

\pandocbounded{\includegraphics[keepaspectratio]{Template_Assignment_files/Template_Assignment_files/figure-latex/demographic_comparison-1.pdf}}

Here you provide a description of why the plot above is relevant to your
specific social problem.

\subsection{3.6 Event Analysis}\label{event-analysis}

\begin{Shaded}
\begin{Highlighting}[]
\FunctionTok{ggplot}\NormalTok{(}\AttributeTok{data =}\NormalTok{ combined\_data, }\FunctionTok{aes}\NormalTok{(}\AttributeTok{x =}\NormalTok{ GDP\_Growth, }\AttributeTok{y =}\NormalTok{ Unemp\_Rate)) }\SpecialCharTok{+}
  \FunctionTok{geom\_point}\NormalTok{(}\AttributeTok{size =} \DecValTok{4}\NormalTok{, }\AttributeTok{color =} \StringTok{"darkblue"}\NormalTok{) }\SpecialCharTok{+}
  \FunctionTok{geom\_text}\NormalTok{(}\FunctionTok{aes}\NormalTok{(}\AttributeTok{label =}\NormalTok{ Year), }\AttributeTok{vjust =} \SpecialCharTok{{-}}\FloatTok{1.2}\NormalTok{, }\AttributeTok{size =} \FloatTok{3.5}\NormalTok{, }\AttributeTok{fontface =} \StringTok{"bold"}\NormalTok{, }\AttributeTok{color =} \StringTok{"black"}\NormalTok{) }\SpecialCharTok{+}
  \FunctionTok{geom\_smooth}\NormalTok{(}\AttributeTok{method =} \StringTok{"lm"}\NormalTok{, }\AttributeTok{color =} \StringTok{"red"}\NormalTok{, }\AttributeTok{linetype =} \StringTok{"dashed"}\NormalTok{, }\AttributeTok{se =} \ConstantTok{FALSE}\NormalTok{) }\SpecialCharTok{+}
  \FunctionTok{theme\_minimal}\NormalTok{() }\SpecialCharTok{+}
  \FunctionTok{labs}\NormalTok{(}
    \AttributeTok{title =} \StringTok{"Event Analysis: Okun\textquotesingle{}s Law and 2020 Policy Distortions"}\NormalTok{,}
    \AttributeTok{subtitle =} \StringTok{"Isolating the 2020 emergency wage{-}subsidy cushion effect"}\NormalTok{,}
    \AttributeTok{x =} \StringTok{"Annual GDP Growth Rate (\%)"}\NormalTok{,}
    \AttributeTok{y =} \StringTok{"Total Unemployment Rate (\%)"}\NormalTok{,}
    \AttributeTok{caption =} \StringTok{"Source: Eurostat Unified Matrix"}
\NormalTok{  ) }\SpecialCharTok{+}
  \FunctionTok{scale\_y\_continuous}\NormalTok{(}\AttributeTok{expand =} \FunctionTok{expansion}\NormalTok{(}\AttributeTok{mult =} \FunctionTok{c}\NormalTok{(}\FloatTok{0.05}\NormalTok{, }\FloatTok{0.12}\NormalTok{)))}
\end{Highlighting}
\end{Shaded}

\begin{verbatim}
## `geom_smooth()` using formula = 'y ~ x'
\end{verbatim}

\pandocbounded{\includegraphics[keepaspectratio]{Template_Assignment_files/Template_Assignment_files/figure-latex/unnamed-chunk-1-1.pdf}}

\section{Part 4 - Discussion}\label{part-4---discussion}

\subsection{4.1 Discuss your findings}\label{discuss-your-findings}

\section{Part 5 - Reproducibility}\label{part-5---reproducibility}

\subsection{5.1 Github repository link}\label{github-repository-link}

Provide the link to your PUBLIC repository here: \ldots{}

\subsection{5.2 Reference list}\label{reference-list}

Use APA referencing throughout your document.

\protect\phantomsection\label{refs}
\begin{CSLReferences}{1}{0}
\bibitem[\citeproctext]{ref-smith2020}
Smith, J. (2020). My article. \emph{Journal of Things}, \emph{1}, 1--10.

\end{CSLReferences}

\end{document}
